% Chapter 1, section 3: The model
% Write in plain LaTeX. Code of a guide step: \guide{<step>}; commands: \begin{shell}...\end{shell};
% a screenshot: \screenshot{file}{caption}; a gap in the draft: \placeholder{...}.
% See book/latex/queen.sty for the full list.

\section{The model}

The model of QueenMQ is easy to understand. It rests on four ideas: partitions, leases,
transactions and a simple interface.

\subsection{Partitions}

The core of the system are the partitions. A QueenMQ partition is a FIFO sub-queue inside a queue,
created on the fly, and a queue is a named collection of those sub-queues. The first push that
names a partition creates it, so there is no partition count to choose and nothing to plan: you
give each entity its own partition, a customer, an order, a conversation, and one queue has held
ten million of them.

\queenfig{kafka-vs-queen}{Hashed partitions share one queue per partition; Queen gives every entity its own.}

QueenMQ partitions are not physical entities like in Kafka: there is no dedicated thread or file
for each partition. A partition is a name, a few rows in the broker's memory and a position for
each consumer group, and its messages live in the log files of its queue.

\why{order matters per entity, not across the whole queue. In Kafka you choose a number of
partitions and hash your keys onto them, so a slow key blocks every key that shares its partition:
the head-of-line blocking of the Smartchat story. With one partition per entity, a slow
conversation waits only for itself.}

\subsection{Leases, not static assignment}

A consumer does not own partitions. When it asks for work, the broker gives it a batch from a
partition that is free, and leases that partition to it for a while (\texttt{leaseTime}). Until
the consumer acks the batch, or the lease runs out, no other consumer of the same group gets that
partition. Then the partition is free again, for whoever asks next.

\queenfig{lease-cycle}{A partition's cycle for one consumer group: free, leased, acked, free again.}

Consumer groups give every reader its own position: each group gets every message, and the
workers inside a group share its work. And reading does not delete: an ack only moves the group's
position, and the messages stay until retention removes them, so a new group can read the history
again.

\why{with no static assignment there is no rebalance. You add a worker and it starts taking work;
a worker dies and its lease runs out, and someone else continues from the same message. The order
of each partition is kept, because a partition has one holder at a time.}

\subsection{Transactions}

A worker's step usually does several things: it finishes the message it read, pushes the next
events, saves some state, and maybe schedules something for later. In QueenMQ that is one call: a
transaction that acks, pushes to any queue, writes to the key-value store and schedules timers, and
commits as a single entry of the log, all or nothing.

\queenfig{one-step-one-commit}{One worker step without and with a transaction.}

\why{without it, every step needs an outbox table, an idempotency table and a cron job, plus the
code that keeps them in step. With it, a crash can never leave a step half done.}

\subsection{A simple interface}

QueenMQ chose a simple API over HTTP: JSON on one port, the same for pushing, consuming,
transactions, KV, timers and locks. The SDKs (JavaScript, Python, Go, Rust, PHP, C++) are thin
layers over it, \texttt{curl} is enough to try everything, and Kafka clients can talk to the
broker through its built-in Kafka port.

\queenfig{thin-clients}{Thin clients; the logic lives in the broker.}

\why{any language with an HTTP client can use QueenMQ, and what you see on the wire is easy to read
and debug.}

\subsection{Wrap up}

In QueenMQ, clients are dumb and the broker is where all the logic lives: the order of each
partition, the leases, the retries and the dead-letter queue, deduplication, timers, state, and
the transactions that tie them together. Your application keeps only its own business.
